El backend del SCDEE se desarrolla en Python 3.12 sobre Django 5.1 y
Django REST Framework 3.15, con PostgreSQL 16 como base de datos
relacional, Redis 7 como \textit{broker} de Celery 5.4 y caché, y MinIO
como almacén de objetos compatible con \acs{s3}. La capa de empaquetado
y orquestación local se gestiona mediante Docker Compose, conforme al
\cref{rnf:deployment}.

El proceso de desarrollo se apoya en un conjunto reducido de herramientas
que cubren cuatro funciones complementarias.

\begin{itemize}

  \item \textbf{Calidad de código.} ruff actúa como \textit{linter} y
  formateador único, sustituyendo a las herramientas tradicionales
  (flake8, isort, black) con un único binario y un único fichero de
  configuración. mypy realiza el análisis estático de tipos sobre todo
  el código de aplicación, con \textit{stubs} específicos para Django y
  \acs{drf}. pre-commit aplica ambas herramientas como \textit{hook} de
  Git antes de cada \textit{commit}.

  \item \textbf{Pruebas.} pytest se utiliza como ejecutor de pruebas
  unitarias y de integración, con pytest-django como integración con el
  \acs{orm} y factory-boy con Faker para construir \textit{fixtures}.
  La cobertura se mide con pytest-cov. Las pruebas de carga se ejecutan
  con Locust contra el \textit{endpoint} \texttt{/api/v1/} para verificar
  el \cref{rnf:concurrency}.

  \item \textbf{Generación automática de documentación y diagramas.} El
  esquema OpenAPI se deriva automáticamente del código mediante
  drf-spectacular (\cref{SEC:DES-API}); los diagramas de clase del
  modelo se generan con \texttt{graph\_models} de django-extensions. Estos artefactos
  se regeneran de forma reproducible y forman parte del material gráfico
  de la memoria.

  \item \textbf{Observabilidad en desarrollo.} El \texttt{docker-compose}
  define un perfil opcional \texttt{monitoring} que arranca pgAdmin
  (PostgreSQL), Redis Commander (Redis), Flower (Celery) y Mailpit
  (captura de \acs{smtp}). Estos servicios no participan del despliegue
  de producción y existen únicamente para inspeccionar el estado del
  sistema durante el desarrollo.

\end{itemize}

Como interfaz unificada de comandos, el proyecto incluye un
\texttt{Makefile} que delega cada operación al \texttt{docker-compose}
subyacente: \texttt{make up} arranca los servicios, \texttt{make check}
ejecuta pruebas, \textit{lint} y comprobación de tipos,
\texttt{make sftp-watcher} lanza el \textit{watcher} de ingesta. Esta
capa de abstracción reduce el coste cognitivo del desarrollo diario y
centraliza la documentación de los comandos disponibles.
